\chapter{Theoretical Background and Prerequisite Learnings}
This chapter discusses the theoretical foundations required for building the supply chain digital twin and agentic framework. It covers graph theory as applied to logistics, pathfinding optimization using Dijkstra's algorithm, the architecture of Large Language Model (LLM) agent systems, and inventory replenishment mathematics. Finally, it outlines the software tools used in this project.

\section{Graph Theory in Supply Chain Networks}
A supply chain can be naturally modeled as a directed graph $G = (V, E)$, where $V$ represents the set of nodes (locations) and $E$ represents the set of directed edges (lanes or shipping routes).
Nodes are classified into tiers based on their role:
\begin{equation}
    V = V_{supplier} \cup V_{chokepoint} \cup V_{port} \cup V_{warehouse} \cup V_{market}
\end{equation}
Directed edges represent the flow of physical goods. Each edge $e = (u, v) \in E$ is associated with attributes that determine the logistics weight. In our digital twin, each edge contains:
\begin{itemize}
    \item \textbf{Baseline Transit Time ($T_{base}$):} The average transit days under normal conditions.
    \item \textbf{Logistics Cost ($C_{base}$):} The freight cost in USD to move cargo across the edge.
    \item \textbf{Friction Penalty ($FP$):} A multiplier representing inherent vulnerability to delays.
\end{itemize}

Under normal conditions, the baseline effective weight is calculated as:
\begin{equation}
    W_{effective} = W_{base} \times (1 + FP)
\end{equation}
By modeling the network as a graph, we can perform complex topological analyses, such as reachability, connectivity, and shortest path calculations, which represent cargo flows.

\section{Pathfinding and Route Optimization}
Route optimization is solved using pathfinding algorithms on the network graph. This project utilizes Dijkstra's algorithm, which finds the shortest path from a single source node to a target node on a weighted graph with non-negative edge weights.

The algorithm maintains a set of visited nodes and a list of tentative distances from the source. At each step, it selects the unvisited node with the smallest tentative distance, updates the distances of its unvisited neighbors (relaxation), and marks it as visited.
For an edge $e = (u, v)$ with weight $w(u, v)$, the relaxation step is defined as:
\begin{equation}
    \text{if } D(u) + w(u, v) < D(v) \text{ then } D(v) = D(u) + w(u, v)
\end{equation}
This project optimizes paths based on two separate objectives:
\begin{enumerate}
    \item \textbf{Transit Time Minimization:} The edge weight $w(u, v)$ is set to the effective transit time in days.
    \item \textbf{Logistics Cost Minimization:} The edge weight $w(u, v)$ is set to the effective logistics cost in USD.
\end{enumerate}
When a shock is injected, edge weights increase, and Dijkstra's algorithm is run dynamically to find the new optimal path, simulating shipping lines rerouting around congested corridors.

\section{LLM-based Multi-Agent Orchestration}
Large Language Models (LLMs) have evolved from text-prediction engines to decision-making agents. An agentic AI framework (such as CrewAI) structure LLMs by assigning roles, goals, and backstories, making them act as autonomous team members.
The key components of an agentic system are:
\begin{itemize}
    \item \textbf{Agents:} Autonomous units initialized with specific LLM configurations, system prompts, and tool access.
    \item \textbf{Tasks:} Concrete assignments given to agents, specifying descriptions, expectations, and output formats.
    \item \textbf{Tools:} Executable functions (e.g., Python functions) that agents can call to interact with the external environment, such as querying databases or updating files.
    \item \textbf{Process Workflow:} The execution policy (sequential, hierarchical, or consensual) that determines how tasks are handed over from one agent to another.
\end{itemize}
In this project, agents use sequential task execution to pass structured messages (JSON payloads) containing shock calculations, routing solutions, cost deltas, and dashboard triggers.

\section{Inventory Replenishment and Stockout Mathematics}
To evaluate the impact of logistics delays, the digital twin incorporates an inventory tracking layer at each consumption market node. The primary metric used is Days of Cover ($DoC$), which represents how long current inventory will last under constant demand:
\begin{equation}
    DoC = \frac{\text{Current Inventory (Metric Tons)}}{\text{Daily Demand Rate (Metric Tons/Day)}}
\end{equation}
When a transit delay ($\Delta_{days}$) occurs on the replenishment lane, the effective lead time increases. If the delay exceeds the days of cover, a stockout occurs:
\begin{equation}
    \text{if } \Delta_{days} > DoC \text{ then Stockout = True}
\end{equation}
To mitigate stockouts, the system triggers emergency restocking. The emergency cost ($C_{emergency}$) is calculated as a penalty proportional to the deficit:
\begin{equation}
    C_{emergency} = \text{Deficit Quantity (MT)} \times \text{Emergency Freight Rate (USD/MT)}
\end{equation}
Additionally, we monitor the Bullwhip Effect, which represents the amplification of demand variability as it moves upstream from markets to warehouses and ports.

\section{Software Tools and Libraries}
The framework is built using the following open-source software tools:
\begin{itemize}
    \item \textbf{Python 3.11:} The core programming language for the backend, optimization, and agent coordination.
    \item \textbf{NetworkX:} A Python library for the creation, manipulation, and study of complex networks.
    \item \textbf{CrewAI:} An orchestration framework for setting up and executing multi-agent workflows.
    \item \textbf{Google Gemini Flash:} The underlying LLM providing linguistic and reasoning capabilities to the agents.
    \item \textbf{FastAPI:} A modern, fast web framework for building RESTful APIs in Python.
    \item \textbf{React \& Vite:} A frontend UI library and build tool for rendering the interactive analytics dashboard.
    \item \textbf{Leaflet:} A JavaScript library for mobile-friendly interactive maps, used to visualize route paths.
\end{itemize}

\section{Summary}
This chapter discussed the theoretical prerequisite concepts, including graph supply chains, Dijkstra's algorithm, agentic AI frameworks, and inventory mathematics. These concepts form the foundation of our 5-layer resilient architecture. The next chapter details the system design, graph schema, and multi-agent setup.
